\section{Appendix V: Speed of Sound, Causality, and EFT Interpretation}
\label{app:eft_causality}

This appendix provides a unified analysis of the physical bounds of the TRXT framework, covering the adiabatic speed of sound, causality constraints, and the domain of validity of the Effective Field Theory (EFT).

\subsection{V.1: Speed of Sound and Causality Analysis}

For a polytropic equation of state $P = K \rho^\gamma$ with $\gamma = 1 + 1/n$, the adiabatic speed of sound $c_s$ is defined as:
\begin{equation}
    c_s^2 \equiv \frac{dP}{d\epsilon} \approx \frac{dP}{d\rho} = K \gamma \rho^{\gamma-1}
\end{equation}

Since $n \approx 1.37$ (from SPARC fits), we have $\gamma \approx 1.73$.
\begin{itemize}
    \item \textbf{High Density Limit:} Since $\gamma > 1$, $c_s$ increases with density.
    \item \textbf{Causality Condition:} The condition $c_s \le 1$ implies a strict density bound:
    \begin{equation}
        \rho \le \rho_{crit} \equiv \left( \frac{1}{K \gamma} \right)^{\frac{1}{\gamma-1}}
    \end{equation}
\end{itemize}
Physical validation requires checking that galactic cores satisfy $\rho_{core} \ll \rho_{crit}$. Given the typical values for SPARC galaxies, $c_s \ll 10^{-3} c$, ensuring strict causality throughout the validated regime.

\subsection{V.2: Range of EFT Validity and the Crossover Scale}
The semi-classical Effective Field Theory (EFT) and the Heat Kernel expansion (Section 2) are strictly valid only in the long-wavelength limit. We define the \textbf{Crossover Scale} $\xi_{cross}$ as the correlation length of the superfluid condensate:
\begin{equation}
    \xi_{cross} \sim \frac{1}{\sqrt{\lambda \langle\Phi\rangle^2}} \approx \ell_{Pl} \times \mathcal{O}(10^3)
\end{equation}
\begin{itemize}
    \item \textbf{IR Regime ($L \gg \xi_{cross}$):} Spacetime appears smooth, and the metric $g_{\mu\nu}$ behaves according to the Einstein-Hilbert action induced via the Heat Kernel. This is the domain of standard General Relativity and the SPARC galaxy dynamics.
    \item \textbf{UV Regime ($L \lesssim \xi_{cross}$):} The semi-classical metric description fails due to the dominance of non-local entanglement. At these scales, the fundamental \textbf{Random Tensor Network} topology and \textbf{TQFT cobordisms} must be used.
\end{itemize}
This crossover explains why GR is an "effective" theory and why the $S \propto A$ area law (Ryu-Takayanagi) is the more fundamental derivation of gravity than the inductive Heat Kernel.

\subsection{V.3: EFT Stability and Operator Classification}

The Nambu-Jona-Lasinio (NJL) Lagrangian employed in this work is understood as an \textbf{effective field theory (EFT)} valid below a UV cutoff $\Lambda_{UV} \sim 0.1 M_{Pl} \approx 10^{18} \text{ GeV}$. Above this scale, the theory requires UV completion.

\textbf{Power Counting and Operator Classification:}
The NJL action has coupling dimension $[G] = -2$, making it power-counting non-renormalizable. Within EFT logic:
\begin{enumerate}
    \item \textbf{Relevant Operators ($[\mathcal{O}] < 4$):} Mass terms — generated dynamically via chiral symmetry breaking.
    \item \textbf{Marginal Operators ($[\mathcal{O}] = 4$):} Kinetic terms — renormalizable.
    \item \textbf{Irrelevant Operators ($[\mathcal{O}] > 4$):} Four-fermion interactions — suppressed by $E^2/\Lambda_{UV}^2$ at low energies.
\end{enumerate}

\subsection{V.4: Asymptotic Safety Perspective}

An alternative UV completion is Weinberg's Asymptotic Safety program, where there exists a non-trivial UV fixed point $G(\mu) \to G^*$. Evidence from functional renormalization group (FRG) studies suggests that gravity + matter systems can exhibit such fixed points, rendering the theory well-defined to arbitrary energies.

\subsection{Conclusion}
The TRXT framework is self-consistent as an EFT below $\Lambda_{UV}$. All falsifiable predictions (SPARC, SIDM, MaVaN) lie well within the EFT validity regime, where corrections are utterly negligible ($\sim 10^{-36}$).
